\section{Limitations}
\label{sec:limits}
Costs were measured on one host, with one model and one serving stack (\Cref{sec:serving}); its disk is
rotational and it has no NVMe and no RDMA fabric, so the remote tier's cost is swept rather than measured,
and the online conclusion holds only if a remote move costs at least about a quarter of a read plus a write
(\Cref{sec:online}). The bill is additive and per block: it has no overlap of moves with computation or
idle link time, no queueing or batching, and no discard-and-recompute action. The only rules we replayed inside a serving
stack are three store-admission rules in vLLM and three write rules in SGLang, on a second host and one trace
(\Cref{sec:contract}), and they show the additive store price wrong in both directions; none of the placement or move results is a statement about latency
under load, and on an asynchronous write path the write half of a move may cost a request nothing, which
the move conclusions do not account for. The replay counts tiers in block slots: at the 64\,MiB block we timed, the GPU tier's 18{,}279 slots
would be a 1.1\,TiB pool, and applying one host's per-block price to a pool of that size assumes that
transfers scale with blocks.
The results use a declared per-tier latency table and are rerun under the measured one; its GPU:CPU
latency ratio is 1:10 where five derived device classes give 1:80 to 1:250, except a deliberately flat one,
a raw copy on our A100 host measures 1:58, and the serving stack, where a GPU-resident block still costs
attention time, about 1:5. Deriving the table from three attention shapes and five device classes puts
the headroom, the share of the heuristic's cost the oracle saves, at 55--98\%, which contains the 92\% the
declared table gives in the same run (93.5\% and 91.8\% on the two traces, \Cref{sec:sources}), but we did
not rerun the ladder or the online replays under the derived tables, and 6 of the 15 configurations
reorder the hierarchy in a way our fixed capacities cannot express.

Every placement result is a single deterministic run with no confidence interval, though the split,
trace length, evaluation window, tier capacities, cost weights, tie rule, action space and billing are swept
rather than fixed; the randomized controls are repeated (twenty request shuffles, ten LRB seeds at the
tightest tier; LeCaR and CACHEUS, also randomized, are single runs), and the ladder, the stationarity test and the age rules are replicated on Alibaba's four traces. The sources
are KV-block hash streams from transformer prefix caching: independent in provenance, not in mechanism,
and we cannot separate workload character from vendor, date or block size. The estimated prior fixes one
pseudo-split, its count bins and its age quartiles; rescaling for its half-prefix counts moves it by at most
0.6 points on the production traces. Our LRB, LeCaR, CACHEUS, S3-FIFO, SIEVE and W-TinyLFU are rebuilt from
their papers, not the authors' artifacts (CACHEUS's CR-LFU expert approximated by our LFU, which keeps counts after eviction, W-TinyLFU's sketch by
exact aged counts), and uniform block sizes remove one of LRB's advantages over LRU, so their results are
statements about this workload rather than about these designs. The best plain rule they are compared
with is chosen per operating point in hindsight; we did not test choosing it from the prefix. Our replays of
these policies bill blocks independently; libCacheSim's implementations of them, counted to a request's first
miss on our traces, move by at most a few points (\Cref{sec:reaction}), and on PrefixSim's own traces that
matching costs ARC at most 1.0 point and S3-FIFO 5.1 at the tightest cache (artifact of \Cref{sec:matching}). Hawkeye and Mockingjay are not run: a KV
block hash carries no analogue of the program counter their predictors are indexed by.
